\section{Shared tools and their degree-corrected specialization}\label{sec:sparse-extension}
The ordinary-SBM article now uses a high-degree-capped adjacency
encoder. Its sharp Bernoulli transfer, same-threshold spectral-row tail,
and deterministic local scalar comparison also apply here. We recall
those tools with their degree-corrected premises. The small-floor
geometry and its exponentially accurate event are proved in full below,
since they differ from the large-threshold adjacency geometry in that
article. The DC classifier in these notes retains its normalized
encoder; no large-threshold DC extension is asserted.

\subsection{Sharp likelihood transfer}\label{sec:sharp}
\begin{lemma}[Bounded tilted error preserves the sharp prefactor]
\label{lem:sharp}
Let $X_j$ be independent Bernoulli variables with probabilities $p_j$,
and let
\[
 T=b_0+\sum_jw_jX_j,\qquad
 K(s)=\log\E e^{-sT}.
\]
Suppose $\abs{w_j}\le W$, the minimizer $s_*$ satisfies
$\tau\le s_*\le1-\tau$ for a fixed $\tau>0$, and under the probability
measure $Q$ proportional to $e^{-s_*T}$,
\begin{equation}\label{eq:sharp-premises}
 \E_Q T=0,\quad vd\le\operatorname{Var}_Q(T)\le V_0d,
 \quad\sum_jr_j\le V_0d,
\end{equation}
where $r_j=Q(X_j=1)$. Put $I=-K(s_*)$.
For every fixed $L<\infty$ and all sufficiently large $d$, uniformly over
$|x|\le L$,
\begin{equation}\label{eq:sharp-oracle}
 c_Ld^{-1/2}e^{-I}\le\PP(T<x)
 \le\PP(T\le x)\le C_Ld^{-1/2}e^{-I}.
\end{equation}

Let $\mathcal W$ be independent of $(X_j)_j$ and let
\[
 \widehat T=T+b+\sum_je_jX_j
\]
have coefficients measurable with respect to $\mathcal W$. On a measurable
training event suppose
\begin{equation}\label{eq:tilt-error}
 \max_j|e_j|\le M,\qquad
 \left|b+\sum_jr_je_j\right|\le B,\qquad
 \sum_jr_je_j^2\le V.
\end{equation}
For fixed $M,V$ and every real $x$,
\begin{equation}\label{eq:sharp-transfer}
 \PP(\widehat T\le x\mid\mathcal W)
 \le C\exp\{s_*x+B+e^M V/2\}\,d^{-1/2}e^{-I}.
\end{equation}
The constant is uniform over the training event. In particular $B$ can
be a nonnegative training-measurable random bound, provided its indicated
exponential moment is integrated when removing the conditioning.
\end{lemma}


\noindent\emph{Proof reference.}
Lemma~A.1 of the SBM article proves this by
Berry--Esseen under the exact Bernoulli tilt, followed by a second tilt
for the independent-training perturbation. That argument makes no
block-model assumption. It therefore applies directly to all scores
used below. See also Zhou and Li~\cite{zhouli} for sharp Chernoff
comparisons in community detection.

\subsection{Application to the actual Bernoulli SBM oracle}

For a true label $a$ and a competitor $b$, the oracle contrast is
\[
 T_{i,ab}=\ell_{i,*}(a)-\ell_{i,*}(b)
 =b_{i,ab}+\sum_{j\ne i}w^{ab}_{z_j}A_{ij},
\]
where
\begin{equation}\label{eq:sbm-weights}
 w_c^{ab}=\log\frac{P_{ac}(1-P_{bc})}{P_{bc}(1-P_{ac})},\qquad
 b_{i,ab}=\sum_{j\ne i}\log\frac{1-P_{az_j}}{1-P_{bz_j}}.
\end{equation}
Its exact Chernoff information is
\begin{equation}\label{eq:sbm-information}
 I_{i,ab}= -\min_{0\le s\le1}
 \sum_{j\ne i}\log\left[
 P_{az_j}^{1-s}P_{bz_j}^{s}
 +(1-P_{az_j})^{1-s}(1-P_{bz_j})^{s}\right].
\end{equation}
On the compact positive full-rank classes used here, probabilities are
comparable to $p$ and bounded away from one, the weights are bounded,
$s_*$ stays away from the endpoints, and the tilted variance is comparable
to $d$. To verify the last assertions, the derivatives of the transform
at its endpoints are respectively minus and plus the two directional
Bernoulli KL divergences, both of order $d$. Its second derivative is
bounded by $Cd$. The row separation implied by
$\sigma_{\min}(P)\ge\kappa p$ gives a lower bound $cd$ on that derivative
throughout $[0,1]$: tilted probabilities are comparable to $p$ and the
sum of squared log-odds differences is at least $cn$.
Thus \cref{lem:sharp} applies uniformly to every pair.

The same conclusion holds for the profile contrasts defined below, with
their own directional information $J_{i,a\to b}$.
It is not necessary for those two directional exponents to coincide.

\subsection{Fixed-\texorpdfstring{$K$}{K} comparison and label alignment}

Let $q_i$ be the error probability of the specified oracle at vertex $i$.
For any coherent family of its candidate scores, strict pairwise loss to a
competitor implies an oracle error; an oracle error implies non-strict
loss to some competitor. Hence
\[
 \max_{b\ne a}\PP(T_{i,ab}<0)
 \le q_i\le\sum_{b\ne a}\PP(T_{i,ab}\le0).
\]
By \eqref{eq:sharp-oracle},
\begin{equation}\label{eq:oracle-pairs}
 q_i\asymp\frac1{\sqrt d}\sum_{b\ne z_i}e^{-I_{i,ab}}
\end{equation}
for the SBM oracle; replace $I$ by $J$ for the profile oracle. The constants
may include a factor depending on fixed $K$.
Fitted partitions are aligned with the true groups by a single common
permutation on the good event. Uniform score comparisons under that
alignment imply
\[
 \E[\ell(\widehat z,z)\ind{\mathcal G}]
 \le\frac1n\sum_i\PP(\text{aligned fitted error at }i,\mathcal G).
\]
The ensuing pairwise bounds apply to fixed true coefficient vectors; a
data-dependent matching does not become an independence assumption.
Finally,
\begin{equation}\label{eq:bad-payment}
 \E\ell(\widehat z,z)
 \le C R_*+\PP(\mathcal G^c).
\end{equation}
\Cref{sec:probability} proves, for the historical decoders, $\PP(\mathcal G^c)=O(R_*)$ under the stated
model conditions. It is this comparison, rather than merely
$\PP(\mathcal G^c)=o(1)$, that completes each theorem.

\subsection{Empirical-floor geometry for bounded degree parameters}
Throughout this section $D_0=\diag(\Omega\one)$,
$G_i=\max(D_i,a_0\bar D)$, and $\norm v_G^2=v^\top Gv$.
The floor multiplier is fixed and sufficiently small.
\begin{proposition}[Sparse degree-floor geometry]\label{prop:sparse-floor}
Under the fixed structural conditions \eqref{eq:mean-model}--%
\eqref{eq:structural}, suppose $d=np\to\infty$ and $np^2\le C_1$.
For every fixed $r>0$, with probability at least $1-C_rn^{-r}$,
$cd\le G_i$, $cnd\le\sum_iD_i$, and
\begin{align}
 \norm{G^{-1/2}(A-\Omega)G^{-1/2}}_{\op}&\le C_r d^{-1/2},
 &\sum_iG_i&\le C_rnd,\label{eq:floor-noise}\\
 \sum_i\frac{(D_i-(D_0)_{ii})^2}{G_i}
 +\sum_i\frac{(G_i-(D_0)_{ii})^2}{G_i}
 &\le C_rn.\label{eq:floor-degree-energy}
\end{align}
Let $U$ contain the $K$ eigenvectors of $S_G$ with largest absolute
eigenvalues, let $\Pi=UU^\top$, and let
$H=G^{-1/2}\Pi G^{1/2}$. For all sufficiently large $n$, on the same
event the retained eigenvalues have a fixed separation from zero, the
remaining eigenvalues have magnitude at most $C_r/\sqrt d$, and
\begin{equation}\label{eq:floor-block-residual}
 \norm{(I-\Pi)G^{1/2}Z}_F\le C_r\sqrt n.
\end{equation}
For every block-constant $f$ with $\norm f_\infty\le C_f$,
\begin{equation}\label{eq:floor-weighted-errors}
 \norm{(I-H)f}_G\le C_rC_f\sqrt n,\qquad
 \norm{(I-H)f}_2^2\le C_rC_f^2n/d,
 \quad \norm v_G^2=v^\top Gv.
\end{equation}
After alignment, the orthonormal population basis
\[
 V=G^{-1/2}\Theta Z
       (Z^\top\Theta G^{-1}\Theta Z)^{-1/2}O
\]
satisfies $\norm{U-V}_F\le C_r/\sqrt d$ and
$\sqrt{G_i}\norm{V_{i\cdot}}_2\le C_r\sqrt{d/n}$.
The same assertions hold simultaneously for recomputed graphs after
any fixed number of deletions, after increasing the probability
parameter. No lower bound proportional to $\log n$ is imposed.
\end{proposition}


\begin{proof}
The Bernoulli transform of the total number of edges gives
$cnd\le\sum_iD_i\le Cnd$ except with probability $e^{-cnd}$.
For a sufficiently small fixed $a_0$, on this event
\[
 cd\le\tau=a_0\bar D\le Cd,
 \qquad \tau<\min_i(D_0)_{ii}.
\]
Let $d_* = n\max_{i\ne j}\E A_{ij}\asymp d$, and put
$W_{h,ii}=\sqrt{\min(1,2d_*/D_i)}$, with value one at zero degree.
The set $\mathcal S=\{i:D_i>2d_*\}$ has size at most $10n/d_*$
except with probability $e^{-cn}$. Indeed, if
$m=\lfloor10n/d_*\rfloor+1\le n$ vertices exceed the threshold,
some set $T$ of size $m$ satisfies $\sum_{i\in T}D_i>2md_*$.
This is an independent-edge sum with coefficients in $\{1,2\}$
and mean at most $md_*$. At $t=(\log2)/2$, its exponential
transform bounds this probability by
$\exp\{-(\log2-1/2)md_*\}$, since
$(e^{tc}-1)/c\le(e^{2t}-1)/2=1/2$ for $c=1,2$.
The union over $T$ costs at most
$\exp\{m\log(en/m)\}=\exp\{o(md_*)\}$.
If $m>n$, the assertion is automatic.

Every row of $W_hAW_h$ has squared Euclidean norm at most
$W_{h,ii}^2D_i\le2d_*$. The regularized-adjacency theorem of
Le, Levina and Vershynin
\cite[Theorem 2.1 and Remark 2.2]{llv2017}, in its row-norm form,
therefore gives
\[
 \norm{W_hAW_h-\E A}_{\op}\le C_r\sqrt d
\]
with failure probability at most $C_rn^{-r}$. It permits these
observed-graph-dependent weights. The difference
$W_h\Omega W_h-\Omega$ is supported on rows or columns in
$\mathcal S$, and every entry has magnitude at most $Cp$; hence
\[
 \norm{W_h\Omega W_h-\Omega}_{\op}
 \le Cp\sqrt{n|\mathcal S|}\le C\sqrt d.
\]
The completed-mean diagonal costs $O(p)$. Thus
$\norm{W_h(A-\Omega)W_h}_{\op}\le C_r\sqrt d$.
The exact factorization
\[
 G^{-1/2}=(2d_*)^{-1/2}FW_h,
 \qquad F_{ii}=\sqrt{\frac{\max(D_i,2d_*)}{G_i}}
\]
has uniformly bounded $F$, because $G_i\ge cd$ and $G_i\ge D_i$.
Conjugating proves the noise bound in \eqref{eq:floor-noise}.
Also $G_i\le D_i+\tau$ gives $\sum_iG_i\le Cnd$.

Applying the noise bound to $G^{1/2}\one$ gives
\[
 \sum_i\frac{(D_i-(D_0)_{ii})^2}{G_i}
 =\norm{G^{-1/2}(A-\Omega)\one}_2^2\le C_rn.
\]
Since $\tau<(D_0)_{ii}$, the maximum defining $G_i$ only reduces
the coordinatewise distance to $(D_0)_{ii}$. This proves the
second energy bound in \eqref{eq:floor-degree-energy}.

Let $J=Z^\top\Theta G^{-1}\Theta Z$. In each block,
Cauchy--Schwarz, the bounded degree parameters, and comparable
block masses give
\[
 \frac{cn}{d}
 \le\frac{(\sum_{i:z_i=a}\theta_i)^2}
          {\sum_{i:z_i=a}G_i}
 \le J_{aa}=\sum_{i:z_i=a}\frac{\theta_i^2}{G_i}
 \le\frac{Cn}{d}.
\]
The nonzero spectrum of $G^{-1/2}\Omega G^{-1/2}$ is that of
$pJ^{1/2}CJ^{1/2}$, which has a fixed separation from zero.
Subspace perturbation and \eqref{eq:floor-noise} give the claimed
eigenvalue and Frobenius bounds. The formula for $V$ and the lower
bound on $J$ give its weighted row bound.

The identity $(D_0)_{ii}=d\theta_iS_{z_i}$ shows that
$G^{-1/2}D_0Z$ belongs to
$\col(G^{-1/2}\Theta Z)$. Moreover
\[
 \norm{G^{1/2}Z-G^{-1/2}D_0Z}_F^2
 =\sum_i\frac{(G_i-(D_0)_{ii})^2}{G_i}\le C_rn,
 \qquad \norm{G^{1/2}Z}_F^2\le C_rnd.
\]
Combining this approximation with the $O(d^{-1/2})$ projector
error proves \eqref{eq:floor-block-residual}.
The matrix $H$ is the $G$-orthogonal projector corresponding to
$\Pi$. Writing $f=Zb$ then gives the first bound in
\eqref{eq:floor-weighted-errors}; the second follows from $G_i\ge cd$.

Deleting a fixed number of vertices changes each block mass and
each sum of degree parameters by $O(1)$. The preceding argument
uses their comparability, not exact mean-one normalization on a
minor, and applies with its recomputed degrees and threshold.
Increasing the fixed probability parameter before a finite-depth
union bound proves the simultaneous assertion.
\end{proof}

\subsection{Spectral-row tails and local deleted-graph comparison}
Let $\mathcal E_r$ denote the aggregate event just described, including
recomputed versions after at most three deletions. Each deleted graph
also has its own training-measurable event. The same-threshold row-tail
proof of Lemma~C.1 in the adjacency SBM article applies to any fixed
multiplier in $\tau/\bar D$, provided the aggregate premises hold.
Indeed, with $W=\tau^{1/2}G^{-1/2}$ and $\widetilde A=WAW$, the preceding
proposition gives $\norm{W(A-\Omega)W}_{\op}\le C\sqrt d$ and
an informative gap of order $d$. Also
\[
 W_i\le1,\quad W_i^2D_i\le\tau,\quad
 \norm{\widetilde A}_{\op}\le\tau,\quad
 \norm{\Omega_{i\cdot}}_2\le Cd/\sqrt n.
\]
These bounds and the exact threshold change
$\tau^{(-i)}-\tau=a_0(\bar D-2D_i)/(n-1)$ are all the
weight and population inputs to that proof. It uses neither
$\tau\ge\max_i(D_0)_{ii}$ nor equality of the degree parameters.
Thus the small multiplier $a_0$ used here gives
\begin{lemma}[Truncated energy for the degree-floor eigenvectors]
\label{restore:tail-energy-lemma}
There are fixed constants $R_0,C,c>0$ such that, for every fixed
$R\ge R_0$ and all sufficiently large $n$,
\begin{equation}
 \E\!\left[\ind{\mathcal E_r}
       \sum_{j=1}^n\norm{U_{j\cdot}}_2^2
                  \ind{\sqrt n\norm{U_{j\cdot}}_2>R}\right]
 \le CR^2e^{-cd\log R}.
 \label{restore:tail-energy}
\end{equation}
The constants $R_0,C,c$ depend on the fixed structural bounds,
$a_0$, and $r$, but not on $R$ or $n$.  The sufficiently-large-$n$
condition may depend on the chosen fixed $R$.  The same result holds
for each recomputed deleted graph with its training-measurable
geometry event.
\end{lemma}


The proof keeps the empirical-threshold perturbation of order $1/n$,
uses unrotated training coefficients in its conditional Bennett bound,
and absorbs the recurrence at the same truncation threshold. It therefore
also applies on the stronger aggregate event below. For an all-zero
training graph its proof-only convention is a deterministic orthonormal
fallback and identity weights; no conditioning on a full-graph event
is introduced.

\begin{corollary}[Exceptional coefficients are hit with exponentially small mass]
\label{restore:tail-hit-corollary}
On a training graph with the preceding floor, put
$H=G^{-1/2}UU^\top G^{1/2}$.
Let $f$ be any training-measurable vector satisfying
$\norm f_\infty\le F_0$, and let $e=(H-I)f$.
For every prescribed fixed $H_0>0$, there is a fixed $M<\infty$
such that, for any independent tested Bernoulli row with probabilities
$p_j\le Cd/n$,
\begin{equation}
 \E\!\left[\ind{\mathcal E_r}
                 \sum_{j:|e_j|>M}p_j\right]
 \le C_{H_0}e^{-H_0d}.
 \label{restore:tail-hit-mass}
\end{equation}
Here $\mathcal E_r$ is the training-measurable geometry event.
No exchangeability among degree parameters is required.
\end{corollary}

\begin{proof}
On that event, $G_j\ge a_-d$ and $\sum_jG_j\le C_Gnd$.
Consequently
\[
 |e_j|
 \le F_0+\frac{\norm{U_{j\cdot}}_2}{\sqrt{G_j}}
                  \norm{U^\top G^{1/2}f}_2
 \le F_0+F_0\sqrt{C_G/a_-}\sqrt n\norm{U_{j\cdot}}_2.
\]
For $M_R=F_0+F_0\sqrt{C_G/a_-}R$, the event
$|e_j|>M_R$ implies $\sqrt n\norm{U_{j\cdot}}_2>R$.
Write $H_R(U)=\sum_j\norm{U_{j\cdot}}_2^2
\ind{\sqrt n\norm{U_{j\cdot}}_2>R}$.
Then \cref{restore:tail-energy-lemma} gives
\begin{equation}
 \E\!\left[\ind{\mathcal E_r}
                    \sum_j\ind{|e_j|>M_R}\right]
 \le\frac n{R^2}\E[H_R(U)\ind{\mathcal E_r}]
 \le Cne^{-cd\log R}.
 \label{restore:tail-coefficient-count}
\end{equation}
Multiplication by $Cd/n$ bounds the expected hit mass by
$Cd e^{-cd\log R}$.
Choose a fixed $R$ with $c\log R>H_0+1$ and absorb the factor $d$.
Independence is required only to interpret the conditional probability
that the tested row hits the training exceptional set; no independence
between the training geometry event and its eigenvectors is asserted.
\end{proof}

\begin{corollary}[Average exceptional target-row mass]
\label{restore:tail-target-corollary}
For every fixed degree cutoff $L$ and every fixed $H_0>0$, there is
a fixed $L_0$ such that
\begin{equation}
 \frac1n\sum_i\PP\!\left\{
    \sqrt{G_i}\norm{U_{i\cdot}}_2>L_0\sqrt{d/n},\quad
    D_i\le Ld,\quad \mathcal E_r\right\}
 \le C_{H_0,L}e^{-H_0d}.
 \label{restore:tail-target-mass}
\end{equation}
\end{corollary}

\begin{proof}
On the displayed event, $G_i\le C_Ld$.
Choose $L_0=R\sqrt{C_L}$.
The target-row exception then implies
$\sqrt n\norm{U_{i\cdot}}_2>R$.
Apply \eqref{restore:tail-energy} and divide its corresponding
expected count by $n$.  Choosing fixed $R$ sufficiently large gives
\eqref{restore:tail-target-mass}.
\end{proof}

\subsection{The local scalar comparison}\label{restore:loo}
\begin{lemma}[A local floored scalar comparison]\label{restore:loo-scalar}
Write $A=\left(\begin{smallmatrix}0&x^\top\\x&B\end{smallmatrix}\right)$,
$k=\sum_jx_j\le Ld$, and let $\Delta$ be the diagonal degree floor
recomputed from $B$. Suppose
\[
 cd\le\Delta_j,\qquad \sum_j\Delta_j\le Cnd,
 \qquad S_B=\Delta^{-1/2}B\Delta^{-1/2}
   =U_B\Lambda_BU_B^\top+N,
\]
where $NU_B=0$, $\norm N_{\op}\le C/\sqrt d$, and the $K$ retained
eigenvalues have a fixed gap from zero. Let the full floor be
$G=\diag(g,\Delta+E)$, with
\[
 cd\le g\le Cd,\qquad E=E_\mathrm{loc}+E_\mathrm{glob},\qquad
 0\le E_\mathrm{loc}\le\diag(x),\quad
 \norm{E_\mathrm{glob}}_{\op}\le Cd/n.
\]
Here $d\to\infty$ and $d\le Cn$. Suppose only on hit coordinates that
\begin{equation}\label{restore:loo-hit}
 \norm{\Delta_j^{-1/2}(U_B)_{j\cdot}}_2
       \le R/\sqrt{nd}\quad\text{whenever }x_j=1.
\end{equation}
Let $U,\Lambda$ be the full informative eigensystem of
$G^{-1/2}AG^{-1/2}$, and put
\[
 H_G=G^{-1/2}UU^\top G^{1/2},\qquad
 H_B=\Delta^{-1/2}U_BU_B^\top\Delta^{1/2}.
\]
For every $\norm w_\infty\le W$, uniformly over the stated inputs,
\begin{equation}\label{restore:loo-main}
 (AH_Gw)_i=x^\top H_Bw_{-i}+O_{L,R,W}(1).
\end{equation}
No bound is assumed on unhit rows of $U_B$ or on $\max_j\Delta_j$.
\end{lemma}


\noindent\emph{Proof reference.}
Lemma~D.1 of the SBM article is deterministic
under precisely these premises. Its generalized eigenvector equations
contain the floored top degree $g$, not the raw incident degree $k$.
Consequently its proof applies unchanged here.

\begin{corollary}[Exact-degree and hollow comparison]\label{restore:loo-corrected}
Under the same local premises, the reconstruction
\eqref{restore:reconstruction} satisfies, for
$\bar w_{-i}=(n-1)^{-1}\sum_{j\ne i}w_j$,
\begin{equation}\label{restore:loo-corrected-action}
 (R^h w)_i=x^\top\{H_B(w_{-i}-\bar w_{-i}\one)
                        +\bar w_{-i}\one\}+O_{L,R,W}(1).
\end{equation}
\end{corollary}
The same deterministic proof, including its local diagonal bound,
is given in Corollary~D.2 of the SBM article. Uniformity in $w$ is algebraic; fitted
coefficients are first replaced by true block contrasts in
\cref{restore:fit-score} before conditioning on the deleted graph.

\subsection{Aggregate probabilities over all degree scales}
\begin{lemma}[Exponentially accurate aggregate geometry above a logarithmic scale]
\label{restore:high-aggregate}
Fix any $A_0>0$ and $H>0$. Under the fixed positive compact structural
conditions, if $d\ge A_0\log n$ and $np^2\le C_1$, there is an aggregate
event $\mathcal E_H$ with
\[
 \PP(\mathcal E_H^c)\le C_H e^{-Hd}
\]
on which the empirical-floor geometry, weighted degree-energy bounds,
two-sided total-degree bounds, and true block-total bounds used in
\cref{restore:fit-geometry-one,restore:fit-geometry-two,restore:fit-block-event}
hold for the full graph and every graph obtained by deleting at most three
vertices. The constants may depend on $A_0,H,C_1$ and the structural
bounds. No condition is imposed on the smallest raw degree or the
largest spectral row.
\end{lemma}


\begin{proof}
Bounded degree parameters give independent upper-triangular entries
with probabilities at most $Cd/n$. The centered bounded-entry norm
inequality of Bandeira--van Handel
\cite[Corollary 3.12 and Remark 3.13]{bvh} gives
\[
 \PP\{\norm{A-\Omega}_{\op}>L_H\sqrt d\}
 \le e^{-(H+2)d}
\]
after choosing a sufficiently large fixed $L_H$. Indeed, the row
variance sum is $O(d)$, the entry bound is one, and the prefactor $n$
in the norm-tail inequality is absorbed using $\log n\le d/A_0$.
The completed-mean diagonal has norm $O(p)$ and is absorbed in the
same threshold. No minimum raw degree is imposed.

The Bernoulli transform of the total number of edges gives, with
failure probability at most $Ce^{-cnd}$,
\[
 cnd\le\sum_iD_i\le Cnd.
\]
For each of the fixed $K^2$ true block totals, Bernstein's inequality
for the independent upper-triangular edges, whose coefficients have
absolute value at most two, gives
\[
 \PP\left\{\left|[Z^\top(A-\Omega)Z]_{ab}\right|>Ln\right\}
 \le 2\exp\left\{-\frac{cL^2n}{d+L}\right\}.
\]
The deterministic completed diagonal costs $O(d)$ and is included
by increasing $L$. Since $d\to\infty$ and $d^2/n\le C_1$,
the exponent is at least $c'L^2d/C_1$ for all sufficiently large $n$.
Increasing the fixed $L=L_H$ therefore makes the failure of all block
totals at most $C_He^{-(H+2)d}$. The total-degree failure is also
bounded by $Ce^{-(H+2)d}$ for large $n$.

These three full-graph events imply the required events on every
fixed-depth minor without a degree-wise probability union. To see this,
let $J=[n]\setminus T$ with $|T|\le3$, set $A_J=A_{JJ}$ and
$\Omega_J=\Omega_{JJ}$, and recompute all degrees and the empirical
floor on $J$. The centered matrix $A_J-\Omega_J$ is a principal
submatrix of $A-\Omega$, so its norm is at most $L_H\sqrt d$.
Removing $T$ changes the total raw degree by at most $2|T|n$;
as $d\to\infty$, the two-sided bound by fixed multiples of $nd$
therefore persists uniformly over $T$. Removing $T$ changes a true
block total of $A-\Omega$ by $O(|T|n)$ deterministically, since
the entries have bounded absolute value. Its $O_H(n)$ bound persists
as well. The surviving group masses and sums of degree parameters differ from
their original values by $O(1)$ and remain comparable to $n$.

The recomputed average degree is consequently comparable to $d$.
For the fixed sufficiently small floor multiplier $a_0$, its threshold
$\tau_J=a_0\overline D_J$ lies between positive fixed multiples of
$d$ and below every coordinate of $D_{0,J}=\Omega_J\one$.
Thus $G_{J,j}=\max(D_{J,j},\tau_J)\ge cd$ and
$\sum_jG_{J,j}\le Cnd$. The raw norm bound immediately gives
\begin{align*}
 \norm{G_J^{-1/2}(A_J-\Omega_J)G_J^{-1/2}}_{\op}
   &\le C_Hd^{-1/2},\\
 \sum_{j\in J}\frac{(D_{J,j}-D_{0,J,j})^2}{G_{J,j}}
   &=\norm{G_J^{-1/2}(A_J-\Omega_J)\one}_2^2
     \le C_Hn.
\end{align*}
Since the population floor is inactive,
$|G_{J,j}-D_{0,J,j}|\le|D_{J,j}-D_{0,J,j}|$; this proves the
other weighted energy bound. The fixed signal gap and the residual
$\norm{(I-U_JU_J^\top)G_J^{1/2}Z_J}_F\le C_H\sqrt n$
now follow by the deterministic Gram-matrix and projection arguments
in \cref{prop:sparse-floor}. Those arguments use only the floor,
total degree, normalized noise, and degree energy, all just proved.
Intersecting the stated aggregate consequences defines
$\mathcal E_H$. The three full-graph events imply it simultaneously
for every indicated minor, proving its claimed failure bound.

For a deleted graph used as independent training data, form the
analogous event using that graph and its further minors. The same
argument gives failure probability $C_He^{-Hd}$, and this event is
measurable from the training graph alone. The full-graph common event
is used only for deterministic implications, never as a condition
that would preserve independence of an incident row.
\end{proof}


\subsection{The degree-corrected tilted cancellation}

Work on a deleted graph so that all coordinates below index actual
neighbors of a separate tested vertex. For the bounded block-constant
linear candidate contrast used in the DC profile reduction,
the exact tilted Bernoulli probabilities have the decomposition
\begin{equation}\label{eq:floor-tilt-decomposition}
 r=D_0h+\epsilon,\qquad
 h\text{ block constant},\quad
 \norm h_\infty\le C/n,\quad \norm\epsilon_\infty\le Cp^2.
\end{equation}
For example the leading term of the exact tilt is
$\theta_i\theta_jpC_{a,z_j}e^{-s w_{z_j}}$, while
$(D_0)_{jj}=\theta_jp\sum_k\theta_k C_{z_j,z_k}$.
Dividing the former by the latter gives the required block-constant
$h_j$. The denominator in the exact Bernoulli tilt contributes
$O(p^2)$ uniformly. This uses the original Bernoulli law throughout.
The target coordinate is absent; artificially appending $r_i=0$ would
instead create an $O(p)$ error at that one coordinate.

For bounded block-constant $f$, put $e=(H-I)f$. The
$G$-orthogonality of $H$ gives the exact cancellation
\begin{align}
 r^\top e
 ={}&-\langle(I-H)h,(I-H)f\rangle_G
       +((D_0-G)h)^\top e+\epsilon^\top e.
       \label{eq:floor-tilted-cancellation}
\end{align}
The first term is $O(1)$ by \eqref{eq:floor-block-residual}, since
the block coefficients of $h$ are $O(1/n)$. For the second,
Cauchy--Schwarz and \eqref{eq:floor-degree-energy} give
\[
 |((D_0-G)h)^\top e|
 \le \norm h_\infty
 \left(\sum_j\frac{((D_0)_{jj}-G_j)^2}{G_j}\right)^{1/2}
 \norm e_G
 \le C.
\]
The last term is at most
$Cp^2\sqrt n\,\sqrt{n/d}\le Cnp^2$. Also $r_j\le Cp$, so
\begin{equation}\label{eq:floor-tilt-moments}
 |r^\top e|\le C(1+np^2),\qquad
 \sum_jr_je_j^2\le Cp\norm e_2^2\le C.
\end{equation}
Thus the first-order bias cancellation and bounded quadratic error do
not require logarithmic degree. The replacement for a uniform
coefficient bound is supplied by
\cref{restore:tail-energy-lemma,restore:tail-hit-corollary}.

\subsection{Two sharp transfers that do not assume every coefficient is bounded}

The next two lemmas give concrete alternatives to a global
row-incoherence event. They concern a row-independent, training-measurable
score, just as in \cref{lem:sharp}.

\begin{lemma}[Exponential-moment coefficient control]\label{lem:sparse-exp-transfer}
Assume the oracle hypotheses of \cref{lem:sharp}. Let
$\widehat T=T+b+\sum_je_jX_j$, with $(b,e)$ measurable from independent
training data. On a measurable training event suppose
\[
 \left|b+\sum_jr_je_j\right|\le B,\qquad
 \sum_jr_je_j^2e^{|e_j|}\le V,
\]
where $B,V$ are fixed. Then, for all real $x$ and all sufficiently
large $d$,
\[
 \PP(\widehat T\le x\mid\text{training})
 \le C e^{s_*x+B+V/2}d^{-1/2}e^{-I}.
\]
The constant is independent of the largest coefficient.
\end{lemma}

\begin{proof}
Taylor's inequality gives, under the original oracle tilt $Q$,
\[
 \log\E e^{-s_*\widehat T}
 \le-I-s_*\left(b+\sum_jr_je_j\right)
       +\frac12\sum_jr_je_j^2e^{|e_j|}
 \le-I+B+V/2.
\]
Tilt again by $e^{-s_*\widehat T}$. The new probabilities satisfy
\[
 \widehat r_j
 =\frac{r_je^{-s_*e_j}}{1-r_j+r_je^{-s_*e_j}}
 \le r_je^{s_*|e_j|}.
\]
For $|e_j|\le1$, this bounds the new centered third absolute moments
by $C_W\sum_jr_j=O(d)$. For $|e_j|>1$, use $s_*\le1-\tau$ and
\[
 e^{s_*|u|}(W+|u|)^3
 \le C_{W,\tau}u^2e^{|u|}\qquad (|u|>1).
\]
The remaining third moments are therefore $O(V)$. The upper variance
bound follows from the same estimates. There is no restriction that
every $e_j$ be bounded.

For the lower variance, restrict to $|e_j|\le1$. The original variance
discarded by this restriction is at most $W^2V$, and the weighted
$\ell_2$ norm of the remaining perturbation is at most $\sqrt V$.
Thus the square root of the original-weight variance of the remaining
$(w_j+e_j)$ is at least $\sqrt{vd-W^2V}-\sqrt V$.
On this restricted set, the new and old Bernoulli variances differ by
only fixed factors. The new total variance is therefore at least $cd$
for large $d$. Berry--Esseen gives a $C/\sqrt d$ bound for every unit
interval under the new tilt. Summing those intervals, using the first
displayed exponential-moment bound, proves the assertion as in
\cref{lem:sharp}.
\end{proof}

\begin{lemma}[Paying exceptional coefficients only when the row hits them]
\label{lem:sparse-truncated-transfer}
Assume the oracle hypotheses of \cref{lem:sharp}, and write $p_j$
for the un-tilted row probabilities. Suppose $p_j\le C_pr_j$, as
holds here for bounded oracle likelihood coefficients. On an independent
training event $\mathcal E$, suppose
\[
 \left|b+\sum_jr_je_j\right|\le B,\qquad \sum_jr_je_j^2\le V.
\]
Here $B,V$ are fixed constants. For any fixed $M>0$ define
\[
 h_M=\sum_{j:|e_j|>M}p_j,\qquad
 \eta_M=\sum_jr_je_j^2\ind{|e_j|>M}.
\]
Then, for every fixed bounded $x$,
\begin{equation}\label{eq:sparse-truncated-transfer}
 \PP(\widehat T\le x)
 \le C_Md^{-1/2}e^{-I}
       +\E[\ind{\mathcal E}h_M]
       +\PP(\mathcal E^c).
\end{equation}
Since $h_M\le C_p\eta_M/M^2$, an expected truncated quadratic error of order
$d^{-1/2}e^{-I}$ suffices. A bound $Ce^{-Hd}$ is a convenient
sufficient condition whenever $I\le C_Id$ and $H>C_I$.
The expected hit-mass condition on $h_M$ is weaker and is itself enough.
\end{lemma}

\begin{proof}
For analysis only, replace $e_j$ by
$e'_j=e_j\ind{|e_j|\le M}$. The truncated coefficients satisfy
\[
 \left|b+\sum_jr_je'_j\right|
 \le B+\eta_M/M\le B+V/M,\qquad
 \sum_jr_j(e'_j)^2\le V.
\]
Hence \cref{lem:sharp} bounds their score error by
$C_Md^{-1/2}e^{-I}$ on $\mathcal E$.
The original and truncated random scores are identical whenever
$X_j=0$ for every $j$ with $|e_j|>M$. Conditional on training, the
probability that this fails is at most
\[
 \sum_{j:|e_j|>M}p_j
 =h_M\le \frac{C_p}{M^2}\eta_M.
\]
Integrate this bound on $\mathcal E$ and pay at most
$\PP(\mathcal E^c)$ outside it. No coefficient truncation is an
additional input or computation required of the implemented decoder.
\end{proof}

For the degree-floor projection there is a concrete reduction to a
weighted row count. Define
\[
 \zeta_j=\sqrt{\frac{nd}{G_j}}\norm{U_{j\cdot}}_2.
\]
If $e=(H-I)f$, $\norm f_\infty\le C_f$ and
$\sum_jG_j\le C_Gnd$, then deterministically
\[
 |e_j|
 \le \frac{\norm{U_{j\cdot}}_2}{\sqrt{G_j}}
        \norm{U^\top G^{1/2}f}_2+C_f
 \le C_f\sqrt{C_G}\,\zeta_j+C_f.
\]
Choosing $M=C_f\sqrt{C_G}L+C_f$ thus gives
$\{|e_j|>M\}\subseteq\{\zeta_j>L\}$. Since $p_j\le Cd/n$,
the sufficient expected-count estimate
\begin{equation}\label{eq:sparse-row-count-target}
 \E\!\left[\ind{\mathcal E}
       \sum_j\ind{\zeta_j>L}\right]\le Cne^{-Hd},
 \qquad H>C_I,
\end{equation}
would imply
$\E[\ind{\mathcal E}h_M]\le Cd e^{-Hd}
=o(d^{-1/2}e^{-I})$ whenever $I\le C_Id$.
Equation \eqref{eq:sparse-row-count-target} is not a consequence of
aggregate geometry alone. It is proved by the same-threshold
truncated-energy argument in \cref{restore:tail-energy-lemma}
and the floor lower bound.

